\subsection*{12.2 Closed Subspace and Projection}

We define the notion of closed subspace below.

\begin{defbox}{12.4}
\end{defbox}

A subset W ofL2(P) is said to be:

A subspace ifX, Y \in W implies\alphaX + \betaY \in W for all scalars \alpha and. \beta

Closed ifXn \in W ,f o rn = 1, 2, 3,... , andXn

L2

-\to X, then. X \in W

In any finite-dimensional vector space, such as Rn, any subspace is closed. The

assumption of closedness is important when we are dealing with vector space with

infinite dimension.

\textbf{Example 12.2.1 (An Example of Subspace that Is Not Closed)} \\

In Example 12.1.2,w el e t W denote the subset of infinite sequences .(x1,x 2,x 3,... ) in which only

finitely many components are nonzero. This is a vector space because it is closed under addition

and scalar multiplication.

However, this subspace is not closed because

Xk \coloneqq(1, 1, 1,..., 1,

k ones

0, 0, 0,...,)

converges to an infinite sequence that contains infinitely many 1's, as k \to\infty , which is not in the

subspace.

\textbf{Example 12.2.2 (Simple Random Variables Associated with a Fixed Partition of the} \\

Sample Space)

Let .(\Omega, \mathcal{F},P) be a probability space and let A1, A2,...,A d be a partition of the sample space

\Omega Let W denote the set consisting of random variables that are constant on each Ak ,f o r k =

1, 2,...,d It is clear that W is a linear subspace of L2(P) .

To show that W is closed, consider a Cauchy sequence .(Xn)\infty

n=1 in W For each index

k = 1, 2,...,d , we can write Xn()1Ak () = cnk 1Ak () for some constant cnk The sequence

(cnk)\infty

n=1 is a Cauchy sequence of real numbers and thus is converging to some constant, say k ,a s

n \to\infty .

Define a random variable X in W by.

\sumd

k=1 k 1Ak S i n c eXk \to X inL2(P) , we see that W is

a closed subspace.

\begin{thmbox}{12.4 (Projection Theorem)}
\end{thmbox}

Given a closed subspace W \subseteq L2(P) and a random variable Y \in L2(P) ,

there exists a unique random variable X* \in W such that

Y - X*2 = inf

X\inW

Y - X2.

This theorem says that we can always find a point in W whose distance to Y is

the smallest among all other points in W .

\textit{Proof} (Existence) Suppose infX\inW Y- X= \alpha. Since. \alpha is the largest lower bound

of .{Y- X : X\in W } , we can find a sequence .(\alphan)\infty

n=1 such that \alphan \alpha , such

that for each n , \alphan is the L2 distance between Y and Xn for some random variable

Xn in the subspace W We claim that the sequence of random variables .(Xn)\infty

n=1 is

a Cauchy sequence.

To prove this claim, we make use of the parallelogram law (Exercise 12.1), which

says that for any two random variables U and V in L2(P) ,

U+ V 2 + U- V 2 =2 U2 +2 V2.

Suppose m and n are two positive integers. In the parallelogram law, substitute U by

Y- Xm and V by Y- XnWe haveU+ V= 2 Y- Xm -Xn andU- V= X n -Xm.

The parallelogram law gives

2Y- X m -X n2 + Xn -X m2 =2 Y- X m2 +2 Y- X n2.

Arrange the terms to get

Xn -X m2 =2 Y- X m2 +2 Y- X n2 - 4 Y- (X n +X m)/22.

In the last term, .(Xn +X m)/2 is the "mid-point" between Xn and. Xm and is thus a

random variable in W (because W is a subspace)The distance Y- (Xn +Xm)/2

cannot be less than \alpha by the definition of infimum. Therefore,

Xn -X m2 \leq2 \alpha2

m +2 \alpha2

n - 4 \alpha2.

Given any \epsilon> 0, we choose a positive integer N such that the right-hand side of

the above inequality is less than \epsilon for all m, n\geq N This is possible because the

sequence .(\alphak)\infty

k=1 is converging to \alpha. This completes the proof of the claim that

(Xn)\infty

n=1 is a Cauchy sequence (with respect to the L2 norm).

By the completeness of L2 norm (Theorem 12.3), the sequence . (Xn)\infty

n=1

converges to a random variable. X* inL2(P) That is, there exists a random variable

X* \inL 2(P) such that Xn converges to X* in mean square. Since W is closed, this

random variable. X* must be in W .

Finally, we can conclude that

\alpha= inf

X\inW

Y- X = limn\to\infty Y- X n= Y- X *.

The last step is due to the continuity of the L2 norm. This proves the existence part.

(Uniqueness) Suppose there are two random variables X*

1 and. X*

2 in W such that

Y - X*

1= Y - X*

2= inf

X\inW

Y - X\coloneqq\alpha.

We can apply the parallelogram law with U = Y - X*

1 andV = Y - X*

2 to obtain

X*

1 - X*

22 \leq 2Y - X*

12 + 2Y - X*

22 - 4\alpha2 = 2\alpha2 + 2\alpha2 - 4\alpha2 = 0.

Therefore,X*

1 - X*

2= 0. It can be true only when X*

1 and. X*

2 are equal as. \hfill $\square$

The infimum in Theorem 12.4 can indeed be achieved. The random variable. X* \in

W satisfiesY - X*2 = minX\inW Y - X2. In view of the previous theorem, we

make the following definition.

\begin{defbox}{12.5}
\end{defbox}

(Projection function) Given a closed subspace W inL2(P) and a random variable

Y \in L2(P) , the unique random variable X that minimizes Y - X2 is called

the projection of Y onto W We denote this minimizer as

ProjW (Y) \coloneqqarg min

X\inW

Y - X2.

Many estimation problems can be formulated as projection onto a closed

subspace. The basic linear regression is an example.

\textbf{Example 12.2.3 (Linear Regression as Projection)} \\

Given n data points .(xi,yi),f o r i = 1, 2,...,n , the problem of linear regression is to find

coefficients a and b that minimize the mean-squared error .

\sumn

i=1(yi - axi - b)2T of o r m u l a t e

this problem as a projection, we can consider a finite sample space \Omega ={ 1, 2,...,n } equipped

with the discrete\sigma-field and the uniform distribution P The data are represented by two functions

X, Y : \Omega \to R,d e fi n e db yX(i) = xi andY(i) = yi ,f o ri = 1, 2,...,n Let I denote the all-one

function on. \Omega .

The mean square error can be expressed as

nE[(Y - aX - bI) 2]= nY - aX - bI2

2.

Define a subspace of L2(P) by

W ={ aX + bI : a \inR,b \inR}.

This subspace is closed because it has finite dimension. We can interpret the optimal choice of a

and b as the projection of the random variable Y onto the closed subspace W .
